\documentclass[11pt]{article}
\usepackage[a4paper,
            top=1.5cm, bottom=1.5cm,
            left=1.5cm,  right=1.5cm]{geometry}

\usepackage{setspace}
\usepackage{microtype}
\usepackage[authoryear]{natbib}
\usepackage{amsmath}
\usepackage{amssymb}

\usepackage{booktabs}
\usepackage{array}
\usepackage{tabularx}

\usepackage{graphicx}
\usepackage{float}
\usepackage{caption}
\captionsetup{font=small, labelfont=bf, justification=justified}

\usepackage[colorlinks=true,
            linkcolor=black,
            citecolor=blue,
            urlcolor=black]{hyperref}


\usepackage{footmisc}


\author{Maria Romagnoli and Lisa Chardon--Denizot}
\title{Applied labour economics}
\date{\today}

\begin{document}

\maketitle
\section{Introduction}
\section{2}
\section{3}
\section{Model}
\label{sec:model}
\subsection{Specification}

Our empirical strategy builds on a stacked difference-in-differences design,
following the approach advocated by \citet{Cengiz2019} and more recently
formalized for multi-period settings by \citet{Baker2022}.\footnote{%
  The stacked design constructs a separate clean "mini-panel" for each
  listing event (2021, 2024, and 2025) defined by a window
  around the relevant reform, and then pools these sub-experiments.
  This avoids the contamination of late-treated units serving as controls
  for early-treated units that plagues two-way fixed-effects estimators
  in heterogeneous-treatment settings \citep{Callaway2021, Sun2021}.}
For practical estimation, rather than a triple-differences specification,
we estimate two separate difference-in-differences regressions
and recover the quantity of interest by comparing the resulting coefficients using the linearity of
conditional expectation.

For each listing cohort $c \in \{2021, 2024, 2025\}$, we estimate the
following stacked event-study regression:

\begin{equation}
  \ln w_{ijrts} =
    \sum_{g \neq -1} \beta_k \bigl(\mathbf{1}[\text{group}_i = g]
      \times \mathbf{1}[\text{policy}_{jrt} = k]\bigr)
    + \mathbf{X}_{ijrt}'\,\gamma
    + \delta_{jrt}
    + \varepsilon_{ijrts},
  \label{eq:eventstudy}
\end{equation}

\noindent
where $w_{ijrts}$ is the monthly base wage of worker $i$ in occupation $j$,
region $r$, at time $t$, belonging to stacked sub-experiment $s$.
The variable $\text{policy}_{jrt} = k$ indexes event time relative to the first
listing of occupation $j$ in region $r$, normalized so that $k = -1$ is the
omitted group (French). $\text{group}_i \in
\{\text{non-EEA},\, \text{EEA},\, \text{French}\}$ identifies the
nationality group of the worker, with French nationals serving as the
reference category.

The individual control vector $\mathbf{X}_{ijrt}$ includes sex, diploma
level, and cumulative experience in months. The fixed-effect structure
$\delta_{jrt}$ absorbs occupation-by-region-by-year variation, which is
the most demanding specification able to control for time-varying labour
scarcity at the occupation-region level, i.e. the shortage premium
effect. Because this high-dimensional fixed
effect is mechanically collinear with the listing indicator, we repeat the
analysis with coarser fixed-effect sets as robustness checks.

The regression is estimated separately for each listing cohort using the
\texttt{feols} function from the \texttt{fixest} package in \textsf{R}
\citep{Berge2018}. Standard errors are clustered at the
occupation-by-region level to account for within-cluster correlation in
both treatment assignment and wages.

\subsection{Parameter of interest and identification}

Let $\hat\beta_k^{\text{non-EEA}}$ and $\hat\beta_k^{\text{EEA}}$ denote
the estimated event-study coefficients for non-EEA and EEA workers
respectively, each relative to French nationals. The parameter of interest
at event time $k$ is the differential post-listing wage effect:

\begin{equation}
  \Delta_k =
  \hat\beta_k^{\text{non-EEA}} - \hat\beta_k^{\text{EEA}}.
  \label{eq:delta}
\end{equation}

Under our hypothesis, we expect $\Delta_k < 0$ for $k \geq 0$:
non-EEA workers should experience a relative wage markdown upon listing,
once the EEA coefficient absorbs any shortage premium accruing to both
groups. The statistical significance of $\Delta_k$ is assessed via the
delta method, which yields an asymptotically valid variance estimate for
linear combinations of estimated coefficients.

Identification rests on four assumptions.
The primary concern is parallel pre-trends: absent listing, non-EEA and
EEA wage trajectories should be parallel to those of French nationals
within the same occupation-region cell. This assumption is testable and
forms the main diagnostic reported below. A second threat is endogenous
worker selection into listed occupations following the reform; to the
extent that selection is positive on unobservables, any estimated markdown
is a lower bound on the true monopsony effect. Third, we assume no general
equilibrium spillovers from listed to unlisted occupations at the
occupation-region level. Fourth, we treat listing events as locally
exogenous, conditional on fixed effects.


\section{Results}
\label{sec:results}

\subsection{Event-study estimates}

Figure %Include event studies graphs%
display the event-study
graphs for the 2021, 2024, and 2025 listing cohorts. Each figure presents
three panels: the estimated $\hat\beta_k^{\text{non-EEA}}$ series (non-EEA
vs.\ French), the estimated $\hat\beta_k^{\text{EEA}}$ series (EEA
vs.\ French), and the implied differential $\Delta_k$ computed via the
delta method. Ninety-five-percent confidence intervals are reported
throughout.

For all three cohorts, the pre-treatment coefficients are statistically
indistinguishable from zero and display no discernible trend, supporting
the parallel pre-trends assumption. This lends credibility to the
identification strategy, even though formal pre-trend tests cannot be
interpreted as proof of the counterfactual assumption.

In the post-treatment periods, neither $\hat\beta_k^{\text{non-EEA}}$ nor
$\hat\beta_k^{\text{EEA}}$ departs significantly from zero for any cohort.
Point estimates are small in magnitude and confidence intervals are wide,
which may reflect genuine heterogeneity in
treatment effects across regions and occupations.

\subsection{Differential effects and delta-method tests}

Table %Include results table
reports the post-listing average of $\Delta_k$ for
each cohort, together with the corresponding standard error and $p$-value
obtained from the delta method.

For the 2021 cohort, no statistically significant differential is detected.
Point estimates of $\Delta_k$ fluctuate around zero in the post-period,
and the null hypothesis $H_0: \Delta_k = 0$ cannot be rejected at any
conventional significance level.  Results for the 2024 cohort are qualitatively similar, with no significant
differential wage effect for non-EEA relative to EEA workers following
listing.

The 2025 cohort yields a weakly significant result at the ten-percent
level for certain post-listing periods. Interestingly, the sign of
$\Delta_k$ is \emph{positive} in this cohort: the point estimate suggests
a slightly higher relative wage for non-EEA workers following listing,
compared to their EEA counterparts. This finding is, if anything,
inconsistent with the monopsony hypothesis; it is more plausibly
attributable to sampling noise given the short post-treatment window
available in the most recent vintage of ForCE.

\subsection{Robustness}

Coarser fixed-effect specifications (e.g., replacing occupation-by-region-by-year
with occupation-by-year fixed effects) do not materially alter the conclusions.

\section{Conclusion}
\label{sec:conclusion}

This paper investigates whether French shortage-occupation lists create
rents that allow firms to depress the wages of non-EEA workers,
whose legal status is tied to remaining in a listed occupation. We embed
this question in a Nash bargaining framework predicting a wage differential
between non-EEA and EEA or French workers following listing, and test it
using a stacked event-study design on ForCE administrative data covering
the 2021, 2024, and 2025 iterations of the policy.

Our main finding is a null result: we do not detect a statistically
significant wage markdown for non-EEA workers relative to EEA workers
following occupation listing, for any of the three reform cohorts
examined. Pre-treatment parallel trends hold throughout, suggesting that
the empirical design is internally valid. The absence of a significant
post-treatment differential is therefore not a mechanical artefact of
the estimation strategy.

Several interpretations of this null are plausible. First, firms may
not exploit workers' permit-conditioned immobility in the short run,
either because wage-setting in these occupations is governed by
collective agreements that limit bilateral bargaining, or because
reputational concerns discourage overt wage discrimination. Second,
non-EEA workers in shortage occupations may self-select on unobservable
productivity that offsets the monopsony markdown, generating an upward
bias that masks a negative structural effect. Third, the window available
for the most recent cohorts (in particular 2025) is too short to
detect dynamic wage effects that may materialise over multiple years.

Our results speak to a growing literature on institutionally induced
monopsony \citep{Naidu2016, Amior2024}. Unlike the UAE kafala setting
studied by \citet{Naidu2016}, the French shortage-list mechanism does not
prohibit job mobility outright; it merely raises its costs. Our null
result is consistent with the hypothesis that moderate rather than
absolute mobility restrictions are insufficient to generate detectable
wage markdowns in a high-regulation labour market where sectoral
bargaining compresses wage dispersion.

Future work should extend the observation window, particularly for the
2024 and 2025 cohorts, as longer panels will improve power and allow
examination of the dynamic path of wage effects. Complementing the wage
analysis with evidence on employment margins (whether listing affects
non-EEA hiring rates or contract duration differentially) would provide
a richer picture of how firms respond to the policy. Finally, exploiting
variation in the stringency of local enforcement or in the share of
workers covered by collective agreements could help disentangle the
mechanisms underlying the null wage result.
\bibliography{biblio}
\bibliographystyle{apalike}

\end{document}